\documentclass[11pt]{article}
\usepackage[margin=1in]{geometry}

% Core packages
\usepackage{amsmath,amssymb}
\usepackage{graphicx}
\usepackage{array}
\usepackage{longtable}
\usepackage{seqsplit}
\usepackage[
  hidelinks,
  pdftitle={Observer-Relative Closure Signatures in Existing LLM Audit Artifacts},
  pdfauthor={Aoi Kawasaki},
  pdfsubject={Technical report},
  pdfkeywords={large language models, structural replay, closure signatures, artifact audit, source-facing annotation}
]{hyperref}

% Paragraphs
\setlength{\parindent}{0pt}
\setlength{\parskip}{1\baselineskip}

\newcommand{\code}[1]{\texttt{\detokenize{#1}}}
\newcommand{\breakcode}[1]{\texttt{\seqsplit{\detokenize{#1}}}}

\title{Observer-Relative Closure Signatures\\
in Existing LLM Audit Artifacts:\\
A Bounded Gate12B Report}
\author{Aoi Kawasaki}
\date{May 6, 2026}

\begin{document}
\maketitle

\begin{abstract}
We study closure signatures in existing LLM audit artifacts using Gate12B, a
read-only secondary audit over Gate12A discrete inference transport outputs.
Across four dense-transformer model lines, archive-family surfaces repeatedly
place high-tension candidates on \code{residual_chord=3} and flat candidates on
\code{residual_chord=1|trusted_tree=2}. The signal survives stricter scale
support, motif observer-scope expansion, and the current basis-preserving local
reparameterization check. A selected source-facing queue aligns high-side R
candidates with \code{conflict-following} rows and flat-side M candidates with
\code{support-following} or \code{non-gluing} rows. Transcript and briefing
sensitivities do not reproduce the same clean archive alignment. We present
this as a bounded artifact study, not a universal interpretability law or
model-quality benchmark.
\end{abstract}

\section{Introduction}

This report studies closure signatures in already materialized LLM audit
artifacts. The object of study is not free generation, leaderboard accuracy, or
internal weight-level causality. The object is a set of replay and transport
artifacts emitted by the Gate12A pipeline and read by Gate12B as a secondary
audit surface.

The narrow motivation is practical. Once a replay artifact contains local
objects, transport relations, triangle cycles, and source-facing text surfaces,
it becomes possible to ask whether some closure-defect signatures survive
changes of observer view, coarse scale, and bounded local reparameterization.
Gate12B asks that question without changing the Gate12A artifact semantics.

The main result is archive-family bounded. In the current dense-transformer
artifact set, archive-family Gate12B surfaces repeatedly place high-tension
candidates on \code{residual_chord=3} and flat candidates on
\code{residual_chord=1|trusted_tree=2}. The same queue and source-facing
annotation path does not make transcript and briefing surfaces reproduce the
same clean archive alignment.

\subsection{What We Claim}

\begin{enumerate}
    \item \textbf{Bounded artifact phenomenon.} We report a repeated
    archive-family closure signature in existing Gate12A/Gate12B audit
    artifacts.
    \item \textbf{Archive-family relation signature.} Across four
    dense-transformer model lines, archive-family Gate12B surfaces repeatedly
    place high-tension candidates on \code{residual_chord=3} and flat candidates
    on \code{residual_chord=1|trusted_tree=2}.
    \item \textbf{Sensitivity survival.} The signal survives the recorded
    observer, scale, motif observer-scope, and basis-preserving local
    reparameterization sensitivities.
    \item \textbf{Source-facing alignment.} A selected archive source-facing
    queue aligns high-side R candidates with \code{conflict-following} rows and
    flat-side M candidates with \code{support-following} or \code{non-gluing}
    rows.
    \item \textbf{Non-archive sensitivity.} Transcript and briefing
    sensitivities do not reproduce the same clean archive alignment under the
    same queue and derived annotation path.
\end{enumerate}

\subsection{What We Do NOT Claim}

\begin{enumerate}
    \item \textbf{No universal law.} We do not claim a universal law of LLM
    reasoning.
    \item \textbf{No benchmark ranking.} We do not present Gate12B as a
    model-quality benchmark or correctness classifier.
    \item \textbf{No answer-quality labels.} Source-facing tags describe the
    direction of a queued source row only; they are not answer-quality labels.
    \item \textbf{No weight-level claim.} We make no claim about model weights
    in general or about mechanistic causality inside a model.
    \item \textbf{No physical invariant.} The bounded reparameterization checks
    are artifact-level stability checks, not a physical gauge invariant.
    \item \textbf{No checkpoint claim.} This report is not a checkpoint,
    release, or replacement for Gate12A.
\end{enumerate}

\section{Controlled Artifact Regime}
\label{sec:artifact-regime}

The empirical surface is deliberately narrow. Gate12B reads existing Gate12A
artifacts and emits secondary audit summaries. No model inference is run for
Gate12B. No Gate12A artifact schema, threshold, classification, or source
semantics is changed.

The selected model lines are:

\begin{itemize}
    \item \code{Qwen/Qwen2.5-0.5B}
    \item \code{Qwen/Qwen2.5-3B-Instruct}
    \item \code{meta-llama/Llama-3.2-3B-Instruct}
    \item \code{Qwen/Qwen3-4B}
\end{itemize}

The selected rendering families are:

\begin{itemize}
    \item \code{archive_v1}
    \item \code{transcript_v1}
    \item \code{briefing_v1}
\end{itemize}

Archive is the family in which the repeated closure signature is reported.
Transcript and briefing are sensitivity surfaces used to test whether the same
structural and source-facing pattern is reproduced outside archive.

\section{Gate12B Operational Framework}
\label{sec:framework}

Gate12B is a read-only secondary audit over Gate12A artifacts. The operational
primitive is:

\begin{quote}
\code{observer x scale x bounded local reparameterization}
\end{quote}

The audit reads existing Gate12A registries:

\begin{itemize}
    \item \code{manifest.json}
    \item \code{explicit_triangle_cycle_registry.jsonl}
    \item \code{triangle_holonomy_registry.jsonl}
    \item \code{transport_relation_registry.jsonl}
\end{itemize}

When source-facing evidence is needed, Gate12B also reads the Gate12A triangle
text-surface audit:

\begin{itemize}
    \item \code{triangle_text_surface_joined.jsonl}
    \item text-surface \code{manifest.json}
\end{itemize}

The representative strict motif setting is:

\begin{quote}
\begin{tabular}{ll}
\code{observer_mode_set} & \code{cycle_motif_expansion_v1} \\
\code{top_k} & \code{3} \\
\code{min_observer_support} & \code{3} \\
\code{min_scale_support} & \code{3} \\
\code{per_band_limit} & \code{2} \\
\end{tabular}
\end{quote}

The source-facing annotation layer is derived by exact-anchor and phrase rules:

\begin{itemize}
    \item exact conflict-anchor text in the answer gives \code{conflict-following}
    \item exact support-anchor text in the answer gives \code{support-following}
    \item no-gluing or not-warranted phrasing gives \code{non-gluing}
    \item otherwise the row is \code{ambiguous}
\end{itemize}

This is not human semantic annotation and not an answer-quality evaluator.

\section{Main Archive-Family Evidence}
\label{sec:archive-evidence}

The main relation-signature result is:

\begin{quote}
\begin{tabular}{ll}
archive high-tension side & \code{residual_chord=3} \\
archive flat side & \code{residual_chord=1|trusted_tree=2} \\
\end{tabular}
\end{quote}

Table~\ref{tab:family-comparison} gives the paper-facing summary.

\begin{table}[ht]
\centering
\small
\caption{Family-level closure-signature summary under the representative strict motif setting.}
\label{tab:family-comparison}
\begin{tabular}{|>{\raggedright\arraybackslash}p{0.18\textwidth}|>{\raggedright\arraybackslash}p{0.28\textwidth}|>{\raggedright\arraybackslash}p{0.28\textwidth}|>{\raggedright\arraybackslash}p{0.16\textwidth}|}
\hline
Family & High-side dominant relation signature & Flat-side dominant relation signature & Reading \\
\hline
\breakcode{archive_v1} & \breakcode{residual_chord=3} & \breakcode{residual_chord=1|trusted_tree=2} & clean archive direction \\
\breakcode{transcript_v1} & \breakcode{residual_chord=1|trusted_tree=2} & \breakcode{residual_chord=3} & mostly reverse / mixed \\
\breakcode{briefing_v1} & \breakcode{residual_chord=1|trusted_tree=2} & \breakcode{residual_chord=3} & mostly reverse / mixed \\
\hline
\end{tabular}
\end{table}

The table should be read as a bounded artifact result. It is not a claim that
archive is universally special, nor that transcript or briefing are failed
families.

\section{Sensitivity And Boundary Results}
\label{sec:sensitivity}

The archive signature survives several recorded sensitivities:

\begin{itemize}
    \item strict scale support at \code{top_k=1}, \code{top_k=3}, and
    \code{top_k=5}
    \item motif observer-scope expansion through \code{cycle_motif_expansion_v1}
    \item the current basis-preserving local reparameterization check
\end{itemize}

The default core observer setting is narrower. Under the stricter
observer-support setting, the default core observer set can make candidates
vanish rather than produce the archive relation-signature result. This boundary
is important because it keeps the motif observer expansion from being treated
as a silent default.

The known Qwen3-4B briefing caveat is a threshold-boundary case:

\begin{itemize}
    \item \code{triangle:000358}
    \item band transition: \code{tense -> flat}
    \item residual delta approximately \code{1.6653345369377348e-16}
    \item not an invariant candidate
    \item not a gauge-stable candidate
\end{itemize}

This is recorded as threshold-boundary behavior, not candidate-level
instability.

\section{Source-Facing Evidence}
\label{sec:source-facing}

The selected source-facing queue contains 16 archive rows:

\begin{itemize}
    \item 8 high-tension rows
    \item 8 flat rows
    \item \code{per_band_limit=2}
\end{itemize}

The derived archive annotation result is:

\begin{quote}
\begin{tabular}{ll}
high-side \code{residual_chord=3} & \code{conflict-following} in 8/8 rows \\
flat-side \code{residual_chord=1|trusted_tree=2} & \code{support-following} in 6/8 rows \\
flat-side \code{residual_chord=1|trusted_tree=2} & \code{non-gluing} in 2/8 rows \\
\end{tabular}
\end{quote}

Table~\ref{tab:source-facing} summarizes archive, transcript, and briefing
annotation counts.

\begin{table}[ht]
\centering
\small
\caption{Source-facing derived annotation summary.}
\label{tab:source-facing}
\begin{tabular}{|>{\raggedright\arraybackslash}p{0.18\textwidth}|r|r|r|r|r|>{\raggedright\arraybackslash}p{0.22\textwidth}|}
\hline
Family & Rows & Support & Conflict & Non-gluing & Ambiguous & Reading \\
\hline
\code{archive_v1} & 16 & 6 & 8 & 2 & 0 & clean high-R conflict / flat-M support-or-non-gluing \\
\code{transcript_v1} & 16 & 10 & 6 & 0 & 0 & relation signatures mostly reverse archive \\
\code{briefing_v1} & 16 & 12 & 4 & 0 & 0 & model-conditioned / mixed \\
\hline
\end{tabular}
\end{table}

These annotation tags describe source-facing direction only. They are not
semantic labels and not answer-quality labels.

\section{Non-Archive Sensitivity}
\label{sec:nonarchive}

The non-archive sensitivity is a defense against a simple annotation-rule
artifact explanation. If the exact-anchor and phrase-rule annotation path
forced every family into the archive pattern, transcript and briefing would
also show the same clean high-R conflict / flat-M support-or-non-gluing
alignment.

They do not.

Under the same representative strict motif setting, transcript and briefing
mostly reverse the relation-signature direction:

\begin{quote}
\begin{tabular}{ll}
non-archive high side & \code{residual_chord=1|trusted_tree=2} \\
non-archive flat side & \code{residual_chord=3} \\
\end{tabular}
\end{quote}

This supports a bounded family-specific reading. It does not imply that
transcript or briefing are incorrect, noisy, or scientifically unimportant.

\section{Limitations}
\label{sec:limits}

The main limitations are:

\begin{itemize}
    \item the study uses existing local artifacts
    \item the current source-facing queue uses \code{per_band_limit=2}
    \item source-facing annotation is exact-anchor / phrase-rule derived
    \item the current dense-transformer set has four model lines
    \item the archive result is family-bounded and artifact-bounded
    \item the current reparameterization check is bounded to the evaluated
    basis-preserving local transform
    \item the evidence package is specified separately and is not yet a public
    release bundle
\end{itemize}

These limits are part of the claim. The report is stronger when the claim
remains bounded to the evidence that earned it.

\section{Artifact Manifest And Reproducibility}
\label{sec:manifest}

The evidence package is specified by
\code{229_GATE12B_PAPER_EVIDENCE_PACKAGE_MANIFEST.md}. The package is
manifest-only at this manuscript-skeleton stage; generated \code{runs/}
artifacts are not committed into this paper source directory.

The required evidence sets are:

\begin{itemize}
    \item motif specificity summary
    \item 36 motif Gate12B runs
    \item archive, transcript, and briefing source queue artifacts
    \item archive, transcript, and briefing source-facing annotation artifacts
    \item archive strict/core boundary runs for appendix sensitivity
\end{itemize}

Recorded integrity status in the manifest:

\begin{itemize}
    \item 36 motif Gate12B runs: checksum mismatches 0
    \item 6 source queue / annotation artifacts: checksum mismatches 0
    \item 24 archive strict/core boundary runs: checksum mismatches 0
    \item 12/12 source queue cases: matching Gate12B/text-surface
    \code{source_gate12a_run_id}
\end{itemize}

\section{Conclusion}

Gate12B reports a bounded archive-family closure-signature phenomenon in
existing LLM audit artifacts. The current dense-transformer archive surfaces
show a repeated high-R / flat-M relation signature that survives the recorded
observer, scale, motif observer-scope, and bounded local reparameterization
sensitivities. Selected source-facing rows align with that structural read,
while transcript and briefing sensitivities do not reproduce the same clean
archive alignment.

The result is paper-ready only as a bounded artifact report. It is not a
universal interpretability law, model-quality benchmark, answer-quality
classifier, physical invariant, checkpoint, or release claim.

\appendix

\section{Evidence Package Manifest}
\label{app:evidence-package}

The paper evidence package is defined in:

\begin{quote}
\code{workstream/229_GATE12B_PAPER_EVIDENCE_PACKAGE_MANIFEST.md}
\end{quote}

The appendix should eventually include a compact copy of the package table.
Until a release package is assembled, the workstream manifest is the source of
truth for which local artifacts support each paper table.

\end{document}
